\chapter{How far can a linear filter go?}\label{ch:floor}

The chapters so far end on 5 November. This chapter and the four after it cover what happened next. Everything here is
still about the 3\um\ to 0.5\um\ problem, unless a section says otherwise.

\section{Why ask this}

Every learned method in Chapter~\ref{ch:results} sat between 30\% and 36\% on COMSOL. Is that close to a hard
limit, or is there a lot of room left? To answer that, we first need the best score a simple \emph{linear} filter could
possibly reach. A linear filter cannot invent anything: it can only boost or damp each frequency that is in the data.

\section{The Wiener filter, fitted on COMSOL}

Our old Wiener baseline was $h/(h^2+\mu)$ with one number $\mu=0.0316$, chosen on simulated data. The real Wiener
filter uses how strong the signal and the noise are at each frequency:
\begin{equation}
W(k) = \frac{H(k)\,S_x(k)}{H(k)^2 S_x(k) + S_n(k)},
\end{equation}
where $H(k)=e^{-k\,\Delta z}$ is the blur, $S_x$ is the power of the sharp field and $S_n$ is the power of the noise.
The script \code{scripts/wiener_oracle.py} measures $S_x$ and $S_n$ from the COMSOL pairs themselves. Here ``noise'' means
whatever the blur model cannot explain: COMSOL's mesh speckle plus any mismatch with free-space physics.

It fits the filter in three ways:
\begin{itemize}
\item \textbf{Oracle:} fitted on the same pairs it is scored on. This is not a usable method; it is a ceiling.
\item \textbf{Other coil:} fitted on one coil, scored on the other.
\item \textbf{Other coil and other heights:} the most honest test.
\end{itemize}

\begin{center}
\begin{tabular}{lc}
\toprule
Filter (local data, two coils, 3\um\ to 0.5\um) & Error \\
\midrule
Old Tikhonov, $\mu=0.0316$ from simulated noise & 49.7\% \\
Scalar $\mu$, fitted on the other coil & 37.3\% \\
Real Wiener, fitted on the other coil & 37.8\% \\
Real Wiener, other coil and other heights & 36.1\% \\
Real Wiener, oracle (ceiling for linear filters) & 34.4\% \\
\midrule
FISTA-300, for reference (15 pairs) & 30.5\% \\
\bottomrule
\end{tabular}
\end{center}

\begin{keyidea}[What the Wiener test showed]
\begin{itemize}
\item The old baseline was badly tuned. Its $\mu$ assumed noise about 18 times stronger than COMSOL really has.
\item Almost all the gain comes from that one number. A single, well-chosen $\mu$ does as well as the full
frequency-shaped filter from 3\um\ up.
\item Even the oracle, at 34.4\%, stays behind FISTA-300 and FISTA-Net. They beat every possible linear filter.
\end{itemize}
\end{keyidea}

\begin{careful}
The fitted Wiener filters saw COMSOL data from another coil; the old $\mu$ saw none. Part of their gain is simply
access to real COMSOL statistics, which a method trained only on simulation does not have. These numbers also come from
the two local coils only. The server rerun with data3 has not been done.
\end{careful}

\section{The lost-band floor}

There is a cleaner way to see the limit. At each frequency, compare the signal that survives the blur with the noise.
Above the frequency where the two are equal (signal-to-noise ratio 1), the data carry no usable information. Now take
the \emph{true} target and delete everything above that frequency. Its error is the \textbf{lost-band floor}: what you
lose if you recover every usable frequency perfectly and invent nothing. The script is
\code{scripts/information_budget.py}.

\begin{center}
\small
\begin{tabular}{lccccc}
\toprule
Input height (\textmu m), target 0.5\um & 1 & 2 & 3 & 4 & 5 \\
\midrule
Floor, L1\_C1 & 17.5\% & 21.0\% & 33.4\% & 39.8\% & 41.3\% \\
Floor, IC\_4um & -- & -- & 35.6\% & 42.6\% & 52.4\% \\
Wiener oracle, mean of coils & 18.0\% & 23.8\% & 34.4\% & 41.2\% & 45.9\% \\
\midrule
Input height (\textmu m) & 6 & 7 & 8 & 9 & 10 \\
\midrule
Floor, L1\_C1 & 42.6\% & 44.8\% & 47.4\% & 53.6\% & 56.0\% \\
Floor, IC\_4um & 58.2\% & 65.4\% & -- & -- & 73.5\% \\
Wiener oracle, mean of coils & 51.0\% & 56.7\% & 50.3\% & 55.3\% & 72.8\% \\
\bottomrule
\end{tabular}
\end{center}

Heights 1, 2, 8 and 9\um\ exist for L1\_C1 only, so their oracle row covers one coil. IC\_4um is harder at every height.

At 3\um\ the floor over both coils is 34.5\%, the same as the Wiener oracle. That is no coincidence: the Wiener filter is
close to ``keep the usable band, drop the rest''.

\begin{keyidea}[The most important number in these chapters]
At 3\um\ to 0.5\um, about a third of the target lies in frequencies the data cannot see. Any method that scores below
34.5\% must be \emph{filling in} those frequencies from a prior: an idea of what sources look like. FISTA does this with its
dictionary, and MAGIC-UNet (Chapter~\ref{ch:magic_results}) with what it learned from simulated wires. From 4\um\ up, 40 to
70\% of the target must come from the prior. So at large standoffs the prior matters more than the algorithm.
\end{keyidea}

\begin{defend}[If your professor asks: why are all the scores near 30\%?]
Because a third of the target is not in the measurement at all. The lost-band floor at 3\um\ is 34.5\%, and the best
linear filter reaches 34.4\%. FISTA-300 (30.5\%) and MAGIC-UNet (27.0\%) beat it only because their priors fill in part of
the lost band.
\end{defend}
